%%%%%%%%%%%%%%%%%%%%%%%%%%%%%%%%%%%%%%%%%
% The Legrand Orange Book
% LaTeX Template
% Version 2.1.1 (14/2/16)
%
% This template has been downloaded from:
% http://www.LaTeXTemplates.com
%
% Original author:
% Mathias Legrand (legrand.mathias@gmail.com) with modifications by:
% Vel (vel@latextemplates.com)
%
% License:
% CC BY-NC-SA 3.0 (http://creativecommons.org/licenses/by-nc-sa/3.0/)
%
% Compiling this template:
% This template uses biber for its bibliography and makeindex for its index.
% When you first open the template, compile it from the command line with the 
% commands below to make sure your LaTeX distribution is configured correctly:
%
% 1) pdflatex main
% 2) makeindex main.idx -s StyleInd.ist
% 3) biber main
% 4) pdflatex main x 2
%
% After this, when you wish to update the bibliography/index use the appropriate
% command above and make sure to compile with pdflatex several times 
% afterwards to propagate your changes to the document.
%
% This template also uses a number of packages which may need to be
% updated to the newest versions for the template to compile. It is strongly
% recommended you update your LaTeX distribution if you have any
% compilation errors.
%
% Important note:
% Chapter heading images should have a 2:1 width:height ratio,
% e.g. 920px width and 460px height.
%
%%%%%%%%%%%%%%%%%%%%%%%%%%%%%%%%%%%%%%%%%

%----------------------------------------------------------------------------------------
%	PACKAGES AND OTHER DOCUMENT CONFIGURATIONS
%----------------------------------------------------------------------------------------

\documentclass[11pt,fleqn]{book} % Default font size and left-justified equations

\usepackage{ifthen} 
\newboolean{VersionProf}
\setboolean{VersionProf}{false}

\usepackage{wasysym}
\usepackage{esint} % various fancy integral symbols
\usepackage{tikz}
\usepackage{tikz-3dplot}
\usepackage{pgfplots}
\pgfplotsset{compat=1.17}
\usetikzlibrary{patterns}
\usepackage{tkz-euclide} 
\usepackage{mathtools}
\usepackage{xurl}
\usepackage[version=4,arrows=pgf-filled,
textfontname=sffamily,
mathfontname=mathsf]{mhchem}
\usepackage{gensymb}
\usepackage{tabularray}
\usepackage{csquotes}

\usepackage[french]{babel}
\usepackage{dingbat}

\usepackage{environ}
\usepackage{wrapfig}

\NewEnviron{solution}{%
    \ifthenelse{\boolean{VersionProf}}{\begin{solutionA}
        \BODY%*
        \end{solutionA}
    }{%
        \relax%
    }%
}%

\usetikzlibrary {arrows.meta} 

\newcommand{\degC}{{\degree}C}
\newcommand{\dd}{\mathrm{d}}
\newcommand{\div}{\mathrm{div}}
\newcommand{\rot}{\vv{\mathrm{div}}}
\newcommand{\grad}{\vv{\mathrm{grad}}}

\tikzset{
    support/.style={
        pattern=north east lines,
        draw=none,
        minimum width=0.3,
        minimum height=0.6
    }
    ,>=latex
}

\usepackage[d]{esvect}

%SPECIFIC TO THIS CHAPTER
\usepackage{physics}
\usepackage{ifthen}
\usepackage[outline]{contour} % glow around text
\usetikzlibrary{calc} % for pic
\usetikzlibrary{angles,quotes} % for pic
\usetikzlibrary{patterns}
\tikzset{>=latex} % for LaTeX arrow head
\contourlength{1.2pt}

\colorlet{myred}{red!65!black}
\definecolor{myred}{RGB}{127,0,255}
\tikzstyle{ground}=[preaction={fill,top color=black!10,bottom color=black!5,shading angle=20},
                    fill,pattern=north east lines,draw=none,minimum width=0.3,minimum height=0.6]
\tikzstyle{mass}=[line width=0.6,myred,fill=myred!40,rounded corners=1,
                  top color=myred!40,bottom color=myred!20,shading angle=20]
\tikzstyle{rope}=[myred!30,line width=1.2,line cap=round] %very thick

% FORCES SWITCH
\tikzstyle{force}=[->,myred,ultra thick,line cap=round]
\tikzstyle{Fproj}=[force,myred!60]
\newboolean{showforces}
\setboolean{showforces}{true}
%----------------------------------------------------------------------------------------

\input{structure} % Insert the commands.tex file which contains the majority of the structure behind the template

\begin{document}

%----------------------------------------------------------------------------------------
%	TABLE OF CONTENTS
%----------------------------------------------------------------------------------------

%\usechapterimagefalse % If you don't want to include a chapter image, use this to toggle images off - it can be enabled later with \usechapterimagetrue

\chapterimage{filmB.jpg} % Table of contents heading image

%\pagestyle{empty} % No headers

%\tableofcontents % Print the table of contents itself

%\cleardoublepage % Forces the first chapter to start on an odd page so it's on the right

%\pagestyle{fancy} % Print headers again



%\part{\'{E}lectromagnétisme}

\setcounter{chapter}{11}

\chapter{Introduction à la physique statistique}

On peut constater que les lois de la mécanique (classique ou quantique) sont toujours réversibles, autrement dit la transformation inverse de l'évolution d'un système mécanique est toujours possible. Ces lois, en principe, sont censées s'appliquer à toutes les particules élémentaires de la matière.

Comment, dès lors, peut-on expliquer les lois de la thermodynamique, et en particulier le second principe qui caractérise la possibilité de transformations irréversibles à l'échelle macroscopique ?

A l'aide de la physique statistique, nous allons tenter de mettre en évidence les propriétés macroscopiques de la matière à partir des lois fondamentales de la mécanique :
\begin{definition}[Physique statistique]
	La physique statistique a pour but de décrire les phénomènes physiques observés à l'échelle macroscopique en combinant une modélisation microscopique (classique ou quantique) avec des statistiques réalisées sur un grand nombre de particules.
\end{definition}

La physique statistique a l'avantage de pouvoir être appliquée aussi bien à partir de la mécanique classique que de la mécanique quantique : on peut alors comparer directement les lois physiques macroscopiques issues des deux modélisations. C'est en particulier par la physique statistique que la loi de Wien (et la loi de Rayleigh-Jeans qui la précédait, avec la catastrophe ultraviolette) ont été obtenues.

\section{Facteur de Boltzmann}

\subsection{Introduction : l'atmosphère isotherme}

Nous allons commencer par un exemple simple : l'étude d'un gaz de température uniforme. En particulier, nous allons étudier l'atmosphère terrestre, dans laquelle la température varie entre environ 300\,K à la surface et jusqu'à 220\,K dans certaines couches. Les variations de température étant de l'ordre de 30\,\% de la température moyenne, on pourra, dans un premier modèle très grossier, les négliger : c'est le modèle de l'atmosphère isotherme.

\subsubsection{Analyse qualitative}

On considère ainsi une colonne d'air, située entre la surface de la Terre prise comme origine $z=0$ et l'espace $z \to +\infty$. Qualitativement, on peut remarquer que :
\begin{itemize}
	\item La gravité va avoir tendance à rassembler les molécules de gaz proche de la surface ;
	\item Mais alors, la pression en surface sera très importante et va venir s'opposer à l'ajout de nouvelles molécules.
\end{itemize}
On s'attend ainsi à avoir un équilibre qui se met en place, entre le gradient vertical de la pression d'une part, et la force de gravité d'autre part.

\subsubsection{Equilibre thermodynamique local}

Pour rendre compte quantitativement de ce phénomène, il va falloir prendre en compte le fait que les particules ne sont pas réparties uniformément : on s'attend à avoir une atmosphère plus dense proche de la surface. Il faut donc définir une densité locale, ce qui fait appel (comme en électrostatique et en magnétostatique) à l'utilisation d'un volume élémentaire mésoscopique.

Cependant, la situation ici est légèrement plus complexe. En électrostatique par exemple, on souhaitait simplement définir la densité volumique de charges, et pour cela il suffisait que le volume soit grand devant la taille moyenne des particules. Ici, c'est différent : dans le volume mésoscopique, nous allons avoir besoin de définir la température $T$ et la pression $p$. Or, ces grandeurs ne sont définies que pour un système à l'équilibre thermodynamique. Ceci nous amène à définir l'équilibre thermodynamique local :
\begin{definition}[Equilibre thermodynamique local]
	On peut considérer qu'un système est à l'équilibre thermodynamique local si on peut découper ce système macroscopique en éléments mésoscopiques qui sont chacun à l'équilibre thermodynamique.
	On peut alors définir une température et une pression pour chacun de ces éléments, ce qui permet de décrire l'objet macroscopique par les champs $T(M,t)$ et $p(M,t)$.
\end{definition}

Il faut donc que l'on choisisse un élément mésoscopique de volume suffisamment grand pour que la température puisse y être définie. Pour que l'agitation thermique s'homogénéise, les particules de gaz doivent avoir la place d'avoir des chocs entre elles, ainsi :
\begin{definition}[Volume mésoscopique en thermodynamique]
	En thermodynamique, un volume mésoscopique $\dd \tau$ est un volume contenant un grand nombre de particules donc grand devant la distance entre les particules), mais également grand devant le libre parcours moyen $\lambda_m$ des particules. Il est petit devant la taille caractéristique du système macroscopique.
\end{definition}
où l'on a utilisé la notion de libre parcours moyen :
\begin{definition}[Libre parcours moyen]
	Le libre parcours moyen est la distance moyenne parcourue par une particule entre deux chocs.
\end{definition}

Pour un gaz parfait, à pression atmosphérique et température ambiante, le libre parcours moyen est typiquement de 100\,nm. Ainsi, une taille caractéristique de 10 à 100\,µm est habituellement adaptée pour un volume mésoscopique.

Sous l'hypothèse d'équilibre thermodynamique local, on peut définir une pression et une masse volumique locales dans l'atmosphère, qui vont dépendre de $z$ : $p = p(z) , \mu = \mu(z)$. 

\subsubsection{Loi du nivellement barométrique}

\begin{capex}
	En écrivant l'équilibre d'une tranche mésoscopique d'air entre $z$ et $z+\dd z$, déterminer une équation reliant le gradient vertical de pression et la masse volumique de l'air.
\end{capex}

On obtient la loi de l'hydrostatique :
\begin{proposition}[Loi fondamentale de l'hydrostatique]
	Soit un fluide au repos soumis à la seule force de pesanteur. On note $p$ la pression dans le fluide, $\mu$ sa masse volumique, et $g$ l'accélération de la pesanteur. Alors :
	\begin{equation}
		\frac{\partial p}{\partial z} = - \mu g
	\end{equation}
\end{proposition}

\begin{capex}
	En déduire que la pression décroît exponentiellement avec l'altitude, avec une hauteur caractéristique $H$ que l'on exprimera en fonction de $R,T,g$ et la masse molaire $M$ de l'air.
\end{capex}

Dans le cas de l'atmosphère terrestre, on peut montrer que $H \approx 8$\,km. Ce modèle très simplifié donne le bon ordre de grandeur pour la hauteur de l'atmosphère, typiquement la pression est réduire de moitié à environ 5\,km d'altitude.

\subsubsection{Interprétation probabiliste}

Nous allons maintenant avoir une approche de physique statistique : nous convertissons une loi macroscopique en une loi à l'échelle microscopique. Bien sûr, pour un état macroscopique donné, il est absolument impossible de savoir où se trouve chacune des molécules de gaz et quelle est sa vitesse. Mais nous pouvons interpréter la loi macroscopique en utilisant les probabilités.

Nous allons, en particulier, raisonner en termes énergétique, afin de pouvoir généraliser notre résultat.

\begin{capex}
	Déterminer l'énergie potentielle de pesanteur d'une particule située à la cote $z$. En déduire la fonction densité de probabilité $\mathbb{P}(E)$ qui caractérise la probabilité d'avoir une énergie $E$ pour une molécule donnée.
\end{capex}

Cette loi nous amène à définir la constante de Boltzmann :
\begin{definition}
	Le constante de Boltzmann est le rapport de la constante des gaz parfaits $R$ et du nombre d'avogadro $\mathcal{N}_A$ :
	\begin{equation}
		k_B = \frac{R}{\mathcal{N}_A} = 1.38 \times 10^{-23}\,\mathrm{J}\cdot\mathrm{K}^{-1}
	\end{equation}
\end{definition}

On constate que la fonction densité de probabilité de l'énergie d'une particule décroît en fonction de l'énergie avec une échelle caractéristique $k_B T$. Cette loi s'interprète physiquement :
\begin{itemize}
	\item Lorsque la température est proche de zéro, alors toutes les particules ont une énergie nulle ($E = 0$) : elles se trouvent dans l'état le plus stable.
	\item En revanche, lorsque la température augmente, on a une compétition entre l'agitation thermique - qui excite les particules - et leur tendance à toujours se rapprocher de l'état fondamental. Les particules vont acquérir une énergie non nulle de l'ordre de $k_B T$.
\end{itemize}

On peut calculer l'énergie potentielle moyenne d'une particule :
\begin{align}
	\langle E \rangle &= \int_0^\infty E \mathbb{P}(E) \dd E \\
	&= \int_0^\infty \frac{E}{k_B T} e^{-\frac{E}{k_B T}}  \dd E \\
	&= k_B T \int_0^\infty u e^{-u} \dd u \\
	&= k_B T \times 1! \\
	&= k_B T
\end{align}

On en déduit l'interprétation physique suivante de $k_B T$ :
\begin{proposition}
	A la température $T$, $k_B T$ représente l'ordre de grandeur de l'énergie d'agitation thermique d'une particule.
\end{proposition}

\subsection{Généralisation}

La propriété que nous avons obtenue se généralise en pratique à toute particule sous certaines hypothèses :
\begin{theorem}[Loi de Boltzmann]
	Soit un système à l'équilibre thermodynamique avec un thermostat de température $T$. La propbabilité pour une particule microscopique indépendante d'être dans un état d'énergie $E$ est proportionnelle au facteur de Boltzmann :
	\begin{equation}
		\mathbb{P}(E) \propto \exp\left(-\frac{E}{k_B T}\right) = e^{-\beta E}
	\end{equation}
	où l'on a utilisé la notation courante $\beta = \frac{1}{k_B T}$.
\end{theorem}
\begin{remark}
	Attention ! L'analogie avec l'atmosphère isotherme n'est pas parfaite, au sens où l'on ne parle pas ici d'une fonction densité de probabilité, mais de la probabilité de se trouver dans un état unique (quantique ou classique) d'énergie $E$. La loi de Boltzmann peut se démontrer, mais la démonstration implique de nombreuses notions de physique statistique qui ne sont pas au programme.
\end{remark}

Le facteur de proportionnalité se calcule en général avec le condition de normalisation :
\begin{proposition}[Condition de normalisation]
	La somme des probabilités de se trouver dans chaque état est égale à 1 :
	\begin{equation}
		\sum_i \mathbb{P}(E_i) = 1
	\end{equation}
\end{proposition}

Une conséquence importante de la loi de Boltzmann est que lorsque des états avec des énergies différentes sont accessibles, on va pouvoir quantifier de combien les probabilités diffèrent en fonction de l'écart entre les niveaux d'énergie :
\begin{itemize}
	\item Si $\Delta E \ll k_B T$ alors les états ont des probabilités similaires, autrement dit tous les états sont équiprobables et les particules vont occuper de la même manière tous les états accessibles ;
	\item Si $\Delta E \gg k_B T$ alors seuls les états d'énergie la plus basse seront occupés. L'agitation thermique n'est pas suffisante pour en extraire les particules.
	\item Si $\Delta E \sim k_B T$ alors les états les moins énergétiques seront les plus occupés, mais les états plus énergétiques seront également occupés.
\end{itemize}

\begin{capex}
	Préciser dans lequel de chacun de ces cas on se trouve pour :
	\begin{itemize}
		\item Un gaz dans une bouteille ;
		\item De l'eau dans une bouteille ;
		\item L'air dans l'atmosphère ;
		\item L'électron de l'atome d'hydrogène.
	\end{itemize}
\end{capex}

\section{Système à deux niveaux}

La loi de Boltzmann est très puissante : elle va nous permettre de décrire le comportement statistique de systèmes classiques, mais surtout de systèmes quantiques pour lesquels on sait que l'énergie est quantifiée. Elle est particulièrement adaptée, par exemple, à la description d'un électron dans un atome. Dans la suite, nous allons considérer un exemple très simple : celui d'un système présentant uniquement deux niveaux d'énergie. Malgré sa simplicité, il a de très nombreuses applications en physique, et le modèle du système à deux niveaux est parfois suffisant pour décrire des phénomènes physiques observables avec une grande précision. Il permet par exemple de modéliser :
\begin{itemize}
	\item Le moment magnétique des électrons, qui est quantifié par le nombre quantique magnétique valant +1/2 ou -1/2. Le modèle du système à deux niveaux permet alors de décrire le magnétisme dans les solides ;
	\item Les noyaux d'hydrogène qui ont les mêmes propriétés magnétiques, on peut alors étudier leur réponse à un champ magnétique imposé. C'est ce qui est utilisé lors d'un IRM (Imagerie par Résonance Magnétique) ;
	\item On peut, de manière générale, assimiler un système à plusieurs niveaux à un système à deux niveaux dans le cas où seuls les deux premiers niveaux sont peuplés significativement, ce qui est fréquent lorsque la température est basse.
\end{itemize}

\subsection{Probabilités de présence et populations}

Puisque l'on peut choisir le niveau d'énergie, on supposera que l'énergie de chacun des deux niveaux est donnée par $\pm \epsilon$ : l'état fondamental a une énergie $E_1=-\epsilon$ et l'état excité une énergie $E_2 = +\epsilon$. L'écart entre les deux niveaux est alors $\Delta E = 2 \epsilon$.

\begin{capex}
	En exploitant la loi de Boltzmann, déterminer la probabilité de présence d'un atome dans chacun des deux niveaux, notées $p_1$ et $p_2$.
\end{capex}

On considère désormais un système constitué de $N$ particules à deux niveaux. La population moyenne dans chacun des deux états s'écrit alors :
\begin{align}
	\langle N_1 \rangle &= N p_1 = \dfrac{N \exp\left(\frac{\epsilon}{k_B T}\right)}{2 \cosh\left(\frac{\epsilon}{k_B T}\right)} \\
	\langle N_2 \rangle &= N p_1 = \dfrac{N \exp\left(-\frac{\epsilon}{k_B T}\right)}{2 \cosh\left(\frac{\epsilon}{k_B T}\right)}
\end{align}
On peut alors tracer $p_1$ et $p_2$ en fonction de la température :

On constate, comme attendu, que lorsque $k_B T \ll \epsilon$, seul le niveau fondamental est occupé (l'agitation thermique est insuffisante pour exciter les atomes) alors que lorsque $k_B T \gg \epsilon$, les deux niveaux sont équiprobables.
\begin{remark}
	\textbf{Attention ! } A haute température, les particules ne sont pas toutes dans l'état excité. Cette affirmation est une erreur classique.
	A haute température, on a équiprobabilité entre les niveaux car c'est la configuration qui maximise l'entropie du système, de la même manière qu'un gaz ne va pas rester d'un côté d'un récipient mais va occuper tout l'espace qui est disponible.
\end{remark}

\subsection{Energie moyenne et fluctuations}

On définit, en physique statistique, les grandeurs moyennes à partir de la loi de probabilité qui leur correspond :
\begin{definition}
	Soit $X$ une grandeur caractérisant un système thermodynamique. En physique statistique, on définit la valeur moyenne de $X$ comme son espérance, en utilisant la loi de probabilité calculée (en général) à l'aide de la loi de Boltzmann :
	\begin{equation}
		\langle X \rangle = \mathbb{E}(X) = \sum_{\mathrm{Etats} n} X_i \mathbb{P}(X_i)
	\end{equation}
\end{definition}

\begin{capex}
	Déterminer l'énergie moyenne et son écart-type en fonction de la température $T$ et du nombre de particules $N$ dans le système. Interpréter les résultats à basse et à haute température.
\end{capex}

On remarque quelques propriétés intéressantes, tout d'abord sur l'énergie moyenne du système :
\begin{itemize}
	\item L'énergie moyenne est toujours négative. Ce n'est pas une propriété générale, ceci dépend de la convention adoptée pour l'origine des énergies. Mais ici, ceci indique que \textbf{le niveau fondamental est toujours plus peuplé que le niveau excité} ;
	\item L'énergie est toujours une fonction croissante de la température.
\end{itemize}

Par ailleurs, l'écart type de cette énergie caractérise les fluctuations d'énergie du système :
\begin{definition}[Fluctuations]
	On appelle fluctuation d'une propriété $X$ d'un système en contact avec un thermostat, les variations de cette grandeur en fonction du temps et/ou de l'espace dues à la nature probabiliste de la présence des particules constituant le système dans les niveaux d'énergie.
	L'amplitude caractéristique des fluctuations est mesurée par l'écart-type de la distribution de probabilité de $X$ :
	\begin{equation}
		\Delta X = \sqrt{\left\langle (X - \langle X \rangle)^2 \right \rangle} = \sqrt{\langle X^2 \rangle - \langle X \rangle^2}
	\end{equation}
\end{definition}

On constate alors que les fluctuations décroissent avec $N$ :
\begin{proposition}
	Dans un système thermodynamique constitué de $N$ particules indépendantes en contact avec un thermostat, les fluctuations sont inversement proportionnelles à $\sqrt{N}$.
\end{proposition}

\subsection{Energie interne et capacité thermique}

Dans le cas d'un système macroscopique ($N \gg 1$), les fluctuations d'énergie sont négligeables; On peut alors considérer que l'énergie est toujours égale à $\langle E \rangle$ et est une quantité fixe, ce qui nous mermet de définir l'énergie interne :
\begin{definition}[Energie interne]
	L'énergie interne est définie uniquement pour un système macroscopique ($N \gg 1$). Elle correspond à la moyenne de l'énergie totale de chacune des particules :
	\begin{equation}
		U = \langle E \rangle
	\end{equation}
	$E$ inclut l'énergie cinétique et l'énergie potentielle microscopique des particules.
\end{definition}

On définit alors la capacité thermique du système :
\begin{definition}
	La capacité thermique à volume constant d'un système thermodynamique macroscopique est la dérivée de son énergie interne par rapport à la température, lorsque l'on fait varier la température d'un thermostat avec lequel el système est à l'équilibre, et lorsque le volume du système est maintenu constant :
	\begin{equation}
		C = \frac{\dd U}{\dd T}
	\end{equation}
\end{definition}
\begin{remark}
	Attention, pour un gaz, cette capacité thermique diffère de la capacité thermique à pression constante.
\end{remark}

\begin{capex}
	Déterminer la capacité thermique du système à deux niveaux, et la relier aux fluctuations d'énergie.
\end{capex}

\begin{remark}
	Déterminer la capacité thermique est une étape importante, car il s'agit d'une grandeur mesurable par l'expérience. L'énergie interne n'est pas directement mesurable puisqu'elle est définie à une constante près.
\end{remark}


\section{Capacités thermiques des gaz parfaits}

Nous nous intéressons désormais à l'étude d'un gaz parfait, et nous souhaitons déterminer sa capacité thermique, en s'inspirant du modèle du système à deux niveaux.

Nous cherchons la capacité thermique du gaz à volume constant, ainsi on considère un gaz confiné dans un volume $V$ fixe, qui est constitué de $N$ particules indépendantes (modèle du gaz parfait). Cette configuration est très similaire au modèle du puits de potentiel infini, dans lequel des particules sont confinées selon la direction $x$. Nous allons donc nous reposer sur ce modèle pour démarrer.

\subsection{Système de $N$ particules dans un puits infini à une dimension}


On considère donc un ensemble de $N$ particules dans un puits de potentiel infini de largeur $L$: il peut correspondre, par exemple, à des particules de gaz enfermées dans une bouteille de volume fixé.

On peut remarquer que ce modèle simple permet de traiter également d'autres situations, en pratique toute situation dans laquelle une particule se trouve dans un puits de potentiel (si on accepte l'approximation grossière du puits infini carré), par exemple, l'étude d'électrons dans des atomes ou d'ions dans des solides cristallins. Bien sûr, comme le puits de potentiel n'est pas réellement carré, certains résultats que nous obtiendrons ne sont pas directement transposables à ces systèmes. 

\subsubsection{Probabilités et occupation des niveaux d'énergie}

Nous avons vu dans le cours de mécanique quantique que les niveaux d'énergie de ce système sont donnés par la relation :
\begin{equation}
	E_n = \frac{n^2 h^2}{8 m L^2} \quad ; \quad n \in \mathbb{N}^\ast
\end{equation}
où $m$ est la masse de la particule, $L$ la largeur du puits, $h$ la constante de Planck. 

De la même manière que pour le système à deux niveaux, on peut calculer la probabilité de se trouver dans chacun des états en exploitant la loi de Boltzmann et en utilisant la condition de normalisation :
\begin{equation}
	\forall n \in \mathbb{N}^\ast, \mathbb{P}(E_n) = \frac{1}{Z} \exp\left(-\frac{E_n}{k_B T}\right)
\end{equation}
où $Z$ est la somme de tous les facteurs de Boltzmann, souvent appelée $fonction de partition$ du système :
\begin{equation}
	Z = \sum_{n=1}^\infty \exp\left(-\frac{E_n}{k_B T}\right)
\end{equation}
\begin{remark}
	Ce résultat se généralise à toute particule indépendante, la somme portant alors sur l'ensemble des états du spectre en énergie de la particule. Il est essentiel de savoir retrouver cette formule très rapidement à partir de la loid e Boltzmann et la condition de normalisation.
\end{remark}



\begin{capex}
	Interpréter le rapport des probabilités d'occupation de deux niveaux d'énergie en termes de population. Etudier les cas limites à haute et basse température. A température ambiante, dans quel cas se situe-t-on pour un atome ? Pour un gaz dans une bouteille ?
\end{capex}

Dans le cas particulier du puits infini, cette somme n'est a priori pas calculable analytiquement. 

\subsubsection{Approximation continue}

On peut considérer le cas d'un gaz dans une bouteille, alors, nous avons vu que l'écart entre de xu niveaux d'énergie est très petit devant l'énergie d'agitation thermique :
\begin{equation}
	E_{n+1} - E_n \ll k_B T
\end{equation}
Ceci nous permet d'approximer le spectre discret d'énergie par un spectre continu, en découpant les énergies accessibles en éléments $\dd E$ tels que $\dd E \ll k_B T$ mais $\dd E \gg E_{n+1} - E_n$.

Ce modèle ne revient pas à oublier le spectre quantique de l'atome, on va simplement le remplacer par une quantité continue qui mesure le nombre d'états entre $E$ et $E+\dd E$, et que l'on notera $\dd n(E)$.

\begin{capex}
	A partir de cette approximation, déterminer la forme de la densité de probabilité $\mathbb{P}(E)$ de l'énergie d'une particule indépendante du systèmé. En déduire que l'énergie moyenne d'une particule est donnée par :
	\begin{equation}
		\langle E \rangle = \frac{1}{2} k_B T
	\end{equation}
\end{capex}

On remarque un comportement différent du système à deux niveaux. En effet, ici, la loi de probabilité de l'énergie est auto-similaire, c'est-à-dire que sa forme est toujours la même en fonction de la température. La distribution de probabilités de l'énergie a toujours la même forme, mais elle est resserrée autour des faibles énergies à basse température, et plus étendue à haute température.

\subsection{Théorème d'équipartition de l'énergie}

L'analyse que nous avons faite n'est vraie qu'à une dimension, et nous souhaiterions l'utiliser pour un gaz compris dans un volume à trois dimensions. Pour cela, nous allons généraliser nos résultats à l'aide du théorème d'équirépartition de l'énergie. Les hypothèses du théorème sont similaires à celles que nous avons faites dans la partie précédente. En particulier, nous supposons que les niveaux d'énergie sont très serrés ($E_{n+1}-E_n \ll k_B T$), si bien que l'on peut faire une approximation continue.

\subsubsection{Degré de liberté quadratique}

Nous allons également utiliser les la notion de degré de liberté comme suit (il s'agit d'une définition simplifiée) :
\begin{definition}[Degré de liberté]
	On appelle degré de liberté d'une particule :
	\begin{itemize}
		\item Une coordonnée permettant de positionner la particule dans l'espace ;
		\item La dérivée temporelle d'une telle coordonnée.
	\end{itemize}
\end{definition}

\begin{example}
	En coordonnées cartésiennes, une particule libre de se déplacer dans tout l'espace a six degrés de liberté : $x,y,z$ qui décrivent sa position, et $\dot x,\dot y, \dot z$ qui décrivent sa vitesse.
\end{example}

Le théorème d'équipartirion de l'énergie concerne les degrés de liberté quadratiques :
\begin{definition}[Degré de liberté quadratique]
	Un degré de liberté $\xi_i$ est dit quadratique si l'énergie de la particule est une fonction quadratique de $\xi_i$, c'est-à-dire :
	\begin{equation}
		E = a  \xi_i^2 + b((\xi_j)_{j \neq i})
	\end{equation}
	où $(\xi_j)_{j \neq i}$ représente les éventuels autres degrés de liberté de la particule.
\end{definition}

\begin{capex}
	Etablir le nombre de degrés de liberté quadratiques dans le cas de la particule située dans un puits carré infini unidimensionnel.
\end{capex}

\subsubsection{Théorème d'équipartition}

Nous avons ainsi pu constater que, avec un degré quadratique d'énergie, la contribution à l'énergie moyenne d'une particule est de $\frac{1}{2} k_B T$. 

Ce résultat est général :

\begin{theorem}[Théorème d'équipartition de l'énergie]
	Dans un système à l'équilibre avec un thermostat de température $T$, si les niveaux d'énergie sont suffisament serrés pour que l'on puisse faire l'approximation continue, alors l'énergie moyenne par particule associée à un degré de liberté quadratique est de $\frac{1}{2} k_B T$.
\end{theorem}

\begin{demonstration}

Considérons un degré de liberté $\xi$ pour lequel on peut faire l'approximation continue. On suppose que le nombre d'états $\dd n(\xi) = n(\xi) \dd \xi$ tels que la particule ait son degré de liberté entre $\xi$ et $\xi + \dd \xi$ est constant, autrement dit, il n'y a pas plus d'états quantiques pour une valeur particulière de ce degré de liberté : $n(\xi) = n = cste$. Cette approximation est très largement vérifiée, sauf cas particuliers typiquement quantiques.

Dans ce cas, on peut écrire la fonction densité de probabilité de l'énergie :
\begin{equation}
	\mathbb{P}(E) \propto \exp \left(-\frac{a \xi^2 + b}{k_B T}\right) \propto \exp \left(-\frac{a \xi^2}{k_B T}\right) 
\end{equation}
où l'on a tiré parti du fait que $b$ ne dépend pas de $\xi$. 
On peut alors remarquer que, lorsque les autres degrés de liberté sont fixés, une variation d'énergie entre $E$ et $E+\dd E$ est équivalente à une variation de la valeur absolue du degré de liberté entre $\xi$ et $\xi + \dd \xi$, pour $\xi>0$:
\begin{equation}
	\mathbb{P}(E) \dd E = \mathbb{P}(\xi) \dd \xi + \mathbb{P}(-\xi) |\dd(-\xi)| = 2 \mathbb{P}(\xi) \dd \xi \quad \mathrm{donc} \quad \mathbb{P}(\xi) =\frac{1}{2} \mathbb{P}(E) \frac{\partial E}{\partial \xi}
\end{equation}
Or on sait que $\partial E / \partial \xi = 2 a \xi$ si bien qu'il existe une constante $A$ telle que :
\begin{equation}
	\mathbb{P}(\xi) = A \xi \exp\left(-\frac{a \xi^2}{k_B T}\right) \quad \mathrm{pour} \quad \xi \ge 0
\end{equation}
On peut calculer $A$ en utilisant la normalisation :
\begin{equation}
	\int_{-\infty}^\infty \mathbb{P}(\xi) \dd \xi = 2 \int_0^\infty \mathbb{P}(\xi) \dd \xi = 1
\end{equation}
qui devient alors :
\begin{equation}
	\frac{1}{2A} = \int_0^\infty \xi e^{-\frac{a \xi^2}{k_B T}} \dd \xi = \frac{k_B T}{a} \int_0^\infty u e^{-u^2} \dd u = \frac{k_B T}{a}
\end{equation}
ou encore $A = \frac{a}{2 k_B T}$. On calcule alors la contribution du degré de liberté $\xi$ à l'énergie moyenne :
\begin{align}
	\langle a \xi^2 \rangle &= 2 \int_0^{infty} a \xi^2 \mathbb{P}(\xi) \dd \xi \\
	&= 2 \int_0^{infty} a \xi^2 \frac{a}{2 k_B T} \xi e^{-\frac{a \xi^2}{k_B T}} \dd \xi \\
	&= \frac{1}{2} k_B T \int_0^\infty v e^{-v} \dd v \quad \left( v = \frac{a \xi^2}{k_B T} \right) \\
	&= \frac{1}{2} k_B T
\end{align}
\end{demonstration}

\subsubsection{Application : vitesse quadratique moyenne}

\begin{capex}
	A partir du théorème d'équipartition, retrouver la vitesse quadratique moyenne $\langle {\vv v}^2 \rangle$ des molécules ou atomes dans un gaz parfait.
\end{capex}

\subsection{Application : capacité thermique des gaz parfaits}

Nous avons vu en 1e année que la capacité thermique molaire d'un gaz parfait s'écrit :
\begin{equation}
	C_{v,m} = \frac{nR}{\gamma-1}
\end{equation}
où $\gamma$ est le coefficient de Laplace du gaz. Nous allons ici démontrer cette formule et retrouver la valeur de $\gamma$ pour les gaz parfaits monoatomiques et diatomiques.

\subsubsection{Gaz parfait monoatomique}

\begin{capex}
	En utilisant le théorème d'équipartition de l'énergie, retrouver la formule classique et montrer que $\gamma = \frac{5}{3}$.
\end{capex}

\subsubsection{Gaz parfait diatomique}

Dans le cas d'un gaz parfait diatomique, le nombre de degrés de liberté est distinct. En effet, l'état dans lequel se trouve une particule est caractérisé non seulement par la position de son barycentre, mais également l'orientation de l'axe reliant les deux molécules (degrés de liberté de rotation). Cette orientation est repérée par deux angles ($\theta$ et $\varphi$ en coordonnées sphériques). On a alors plus de degrés de liberté quadratiques :
\begin{itemize}
	\item $x,y,z$ indiquent la position de la particule. Ils ne modifient pas son énergie : ce ne sont pas des degrés de liberté quadratiques.
	\item $\theta, \varphi$, de même, n'influencent pas l'énergie de la particule.
	\item $\dot x, \dot y, \dot z$ sont liés à l'énergie cinétique de translation et sont des degrés de liberté quadratiques :
	\begin{equation}
		E_{c,trans} = \frac{1}{2} m (\dot x^2 + \dot y^2 + \dot z^2)
	\end{equation}
	\item $\dot \theta, \dot \varphi$ sont liés à l'énergie cinétique de rotation de la molécule autour de son axe principal et de la direction perpendiculaire à son axe :
	\begin{equation}
		E_{c,rot} = \frac{1}{2} (J_\theta \dot \theta^2 + J_\varphi \dot \varphi^2)
	\end{equation}
\end{itemize}

On a donc en tout 5 degrés de liberté quadratiques, ainsi, pour une particule :
\begin{equation}
	\langle E \rangle = \frac{5}{2} k_B T
\end{equation}
d'où, pour $n$ moles de gaz :
\begin{equation}
	U = \frac{5}{2} n R T \quad \mathrm{et} \quad C_{v,m} = \frac{5}{2} n R
\end{equation}
on en déduit directement que pour un gaz parfait diatomique :
\begin{equation}
	\gamma = \frac{7}{5}
\end{equation}
C'est notamment le cas de l'air, qui s'assimile bien à un gaz parfait à pression et température ambiante, et a fortiori à toutes les pressions rencontrées dans l'atmosphère qui sont nécessairement plus basses que la pression de surface.

\section{Capacité thermique des solides cristallins}

La capacité thermique d'un solide cristallin provient de la capacité thermique du réseau cristallin lui-même mais peut avoir d'autres contributions :
\begin{itemize}
	\item Dans le cas d'un métal, les électrons se déplacent librement dans le réseau et peuvent être assimilés à un gaz parfait ; ils ont donc une capacité thermique.
	\item Dans le cas d'un solide magnétique, l'interaction entre les moments magnétiques est source d'une énergie potentielle également responsable d'une capacité thermique.
\end{itemize}
Dans le cadre de ce chapitre, nous nous limitons à une situation dans laquelle aucun de ces deux phénomènes n'intervient (par exemple, nous pouvons étudier un solide ionique). Dans ce cas, nous avons uniquement à décrire le réseau cristallin.

\subsection{Modèle d'Einstein}

Dans le modèle d'Einstein, les atomes occupent des positions précises correspondant à leur position d'équilibre, et se trouvent dans un minimum de potentiel. Considérons l'un de ces atomes, et prenons comme origine de l'espace sa position d'équilibre. Il a pour énergie cinétique :
\begin{equation}
	E_c = \frac{1}{2} (\dot x^2 + \dot y^2 + \dot z^2)
\end{equation}
Par ailleurs son énergie potentielle s'écrit, en développant à l'ordre 2 :
\begin{align}
	V(x,y,z) = &V_0 + \frac{\partial V}{\partial x} x + \frac{\partial V}{\partial y} y + \frac{\partial V}{\partial z} z \\
	&+ \frac{\partial^2 V}{\partial x^2} \frac{x^2}{2} + \frac{\partial^2 V}{\partial y^2} \frac{y^2}{2} + \frac{\partial^2 V}{\partial z^2} \frac{z^2}{2} \\
	&+ \frac{\partial^2 V}{\partial x \partial y} xy + \frac{\partial^2 V}{\partial x \partial z} xz + \frac{\partial^2 V}{\partial y \partial z} yz
\end{align}
Or on sait que $V$ a un minimum en $O$ donc toutes les dérivées d'ordre 1 s'annulent, et par ailleurs, on peut utiliser les symétries du problème pour établir que les fonctions $\mathbb{P}(x),\mathbb{P}(y),\mathbb{P}(z)$ doivent être paires. Les degrés de liberté étant indépendants, on a par ailleurs :
\begin{equation}
	\langle xy \rangle = \iint x y \mathbb{P}(x) \mathbb{P}(y) \dd x \dd y = \int_{-\infty}^\infty x \mathbb{P}(x) \dd x \times \int_{-\infty}^\infty y \mathbb{P}(y) \dd y = 0
\end{equation}
et de même les autres termes croisés ont une moyenne nulle. Ainsi :
\begin{equation}
	\langle V \rangle = \frac{1}{2} \left( \frac{\partial^2 V}{\partial x^2} \langle x^2 \rangle + \frac{\partial^2 V}{\partial y^2} \langle y^2 \rangle + \frac{\partial^2 V}{\partial z^2} \langle z^2 \rangle \right)
\end{equation}
Ainsi, les trois degrés de liberté de position sont également des degrés de liberté quadratiques.

\begin{remark}
	Cette forme de l'énergie potentielle est analogue à celle d'un oscillateur harmonique : en fait, dans le modèle d'Einstein, les atomes sont élastiquement liés à leur position d'équilibre avec une force de rappel proportionelle à la distance entre l'atome et sa position d'équilibre.
\end{remark}

\subsection{Loi de Dulong et Petit}

\begin{capex}
	A partie du modèle d'Einstein et du théorème d'équipartition de l'énergie, déterminer la capacité thermique d'un solide cristallin non métallique. Le résultat est connu sous le nom de \textbf{Loi de Dulong et Petit} et est à connaître.
\end{capex}

\begin{remark}
	L'approximation continue est moins bien vérifiée pour les solides que pour les gaz. Nous verrons en TD le modèle d'Einstein quantique, qui permet d'estimer le rôle des effets quantiques sur la capacité thermique des solides.
\end{remark}


\end{document}
